\theory{Theory of statistics}{How can limited, noisy observations support reliable conclusions about a much larger population or an uncertain process?}{Statistical theory gives mathematical principles for study design, data collection, summarization, inference, prediction, and decisions. It connects probability with optimization and utility so that methods can be compared and their uncertainty quantified.}{Experiments, surveys, medicine, engineering, economics, machine learning, public policy, and scientific data analysis.}{Sampling distributions, estimators, likelihoods, and error control turn data into quantified conclusions about a population.}

\paragraph{Ronald Fisher}
Fisher developed likelihood-based estimation, analysis of variance, and principles for randomized experimental design. These methods connected data collection with efficient estimation and made experimental uncertainty part of the mathematical model.

\paragraph{Jerzy Neyman and Egon Pearson}
Neyman and Pearson formalized hypothesis testing through error probabilities and power. Their framework distinguishes the rule used to make a decision from the long-run error guarantees attached to that rule \cite{statistics_ref1,statistics_ref2}.
\end{historylist}
\section*{3. Theory}
\subsection*{Core theory}
\paragraph{The problem and scope}
Data are observations from a population or process, often with measurement error and random variation. A statistical model describes how the data could have been generated. Statistical theory asks which questions can be answered, how much information is needed, and how reliable an answer is \cite{statistics_ref1}.

Its tasks include sampling from finite populations, measuring error, describing distributions, estimating parameters, testing hypotheses, predicting outcomes, comparing groups, and reducing dimension. It also studies what happens when assumptions such as independence or a chosen distribution are wrong \cite{statistics_ref2}.

\paragraph{Design before analysis}
Randomization makes treatment groups comparable and supports probability statements. Experimental design chooses factors, controls, replication, and sample size to estimate treatment effects efficiently. Survey sampling chooses a representative sample while accounting for unequal selection probabilities and nonresponse.

Good design can reduce cost and uncertainty before any formula is applied. Observational studies can estimate associations but usually require stronger assumptions to infer causation. Statistical theory separates the information supplied by the design from conclusions added by a model \cite{statistics_ref3}.

\paragraph{Inference}
An estimator is a rule that turns data into a parameter estimate. Bias, variance, consistency, and efficiency compare estimators. A confidence interval gives a procedure that covers the true parameter at a stated long-run rate; it is not a probability statement about a fixed parameter after the interval is observed.

Hypothesis testing compares a null model with alternatives. A p-value measures how surprising data would be under the null, while power measures the chance of detecting a specified alternative. Bayesian inference combines a prior distribution with a likelihood to produce a posterior distribution. Decision theory chooses actions by expected loss or utility \cite{statistics_ref4}.

\paragraph{Models and robustness}
Regression models describe relationships and prediction. Dimension reduction replaces many variables by a smaller informative representation. Model checking examines residuals, dependence, outliers, and calibration. Robust methods reduce sensitivity to incorrect distributional assumptions.

\section*{4. Important theorems and results}
\begin{enumerate}[leftmargin=*,itemsep=0.35em]
\item \textbf{Law of large numbers.} Sample averages approach population means under suitable conditions.

\item \textbf{Central limit theorem.} Properly standardized sums often approach a normal distribution, explaining many large-sample intervals.

\item \textbf{Neyman--Pearson lemma.} For testing two simple hypotheses at a fixed error level, the likelihood-ratio test is most powerful.

\item \textbf{Gauss--Markov theorem.} In a linear model with standard assumptions, least squares is the best linear unbiased estimator.

\item \textbf{Sufficiency principle.} A sufficient statistic retains all information about a parameter contained in the sample under the model.

\item \textbf{Bayes theorem.} Posterior odds equal prior odds multiplied by the likelihood ratio.
\end{enumerate}
\noindent\emph{Source for these results:} \cite{statistics_ref1}.

\theoryapplications
\theoryconnections{theory_of_statistics}{%
\item Probability describes the random mechanism that could produce the data, and statistics reverses that direction to infer unknown features of the mechanism.
\item Linear algebra makes regression and covariance computable, calculus derives likelihood equations, and optimization chooses parameters that fit the observations.
\item Decision theory adds the loss of acting on an estimate, so a statistically good answer is judged by consequences as well as by closeness to the data \cite{statistics_ref1,statistics_ref4}.
}{%
\item Statistical theory helps an experiment separate a real effect from sampling noise by comparing an estimate with its expected variability.
\item In medicine it turns trial data into risk estimates, in engineering it monitors failure rates, and in machine learning it controls prediction error on unseen data.
\item Surveys and public policy depend on the same connection, but only when sampling, missing data, dependence, and the decision costs are modeled explicitly \cite{statistics_ref3,statistics_ref5}.
}
\section*{Glossary}
\begin{description}[leftmargin=*,style=nextline,itemsep=0.3em]
\item[\textbf{Estimator.}] A rule that converts data into an estimate of an unknown parameter; different estimators can be compared by bias, variance, and efficiency.
\item[\textbf{Bias.}] The difference between the expected value of an estimator and the true parameter value; an unbiased estimator has zero bias on average.
\item[\textbf{Confidence interval.}] A procedure that, if repeated many times, covers the true parameter at a stated long-run rate; it is not a probability statement about a fixed parameter.
\item[\textbf{Hypothesis testing.}] A formal procedure for comparing a null model against alternatives using data, controlled by error probabilities.
\item[\textbf{p-value.}] The probability of obtaining data as extreme as or more extreme than observed, assuming the null hypothesis is true; it measures surprise under the null.
\item[\textbf{Power.}] The probability of rejecting the null hypothesis when a specific alternative is true; it measures the ability to detect a real effect.
\item[\textbf{Bayesian inference.}] A method that combines a prior distribution with a likelihood to produce a posterior distribution for unknown parameters.
\item[\textbf{Sufficiency principle.}] A sufficient statistic retains all information about a parameter contained in the sample under the model.
\item[\textbf{Likelihood.}] A function of parameters given observed data, proportional to the probability of the data under the model.
\item[\textbf{Variance.}] The expected squared deviation from the mean; measures spread of an estimator or data.
\item[\textbf{Likelihood-ratio test.}] A test comparing maximized likelihoods under two models; the Neyman–Pearson lemma shows it is most powerful for simple hypotheses.
\item[\textbf{Standard error.}] The standard deviation of the sampling distribution of a statistic.
\end{description}
\section*{7. Sample Questions}
\emph{Exercise reference:} \cite{statistics_ref1}. The cited edition is the source for the exercise set; consult its Exercises section for the official statement and numbering.
\begin{enumerate}[leftmargin=*,itemsep=0.9em]
\item \textbf{Exercise 1.} Let $X_1, \ldots, X_n$ be i.i.d.\ $\mathrm{Bernoulli}(p)$. Compute the maximum likelihood estimator $\hat{p}$ for $p$. \emph{Solution.}$L(p) = \prod_{i=1}^n p^{x_i}(1-p)^{1-x_i} = p^{s}(1-p)^{n-s}$ where $s = \sum x_i$. Then $\ell(p) = s\ln p + (n-s)\ln(1-p)$. Setting $\frac{d\ell}{dp} = \frac{s}{p} - \frac{n-s}{1-p} = 0$ gives $\hat{p} = \frac{s}{n} = \bar{X}$, the sample mean.

\item \textbf{Exercise 2.} Compute $\mathrm{Var}(\bar{X})$ for i.i.d.\ observations with variance $\sigma^2 = 9$ and sample size $n = 36$. Compare with the variance of a single observation. \emph{Solution.}$\mathrm{Var}(\bar{X}) = \frac{\sigma^2}{n} = \frac{9}{36} = 0.25$. The variance of a single observation is $9$. The standard error is $\sqrt{0.25} = 0.5$, compared to $\sqrt{9} = 3$ for a single observation. Averaging reduces the standard deviation by a factor of $\sqrt{n} = 6$.

\item \textbf{Exercise 3.} For the model $Y = \beta_0 + \beta_1 X + \varepsilon$ with $\varepsilon \sim N(0, \sigma^2)$, the least-squares estimator is $\hat{\beta}_1 = \frac{\sum(X_i - \bar{X})(Y_i - \bar{Y})}{\sum(X_i - \bar{X})^2}$. Show that $\hat{\beta}_1$ is unbiased when $\mathrm{E}[\varepsilon_i] = 0$ and the $X_i$ are fixed. \emph{Solution.}$\mathrm{E}[\hat{\beta}_1] = \frac{\sum(X_i - \bar{X})\mathrm{E}[Y_i - \bar{Y}]}{\sum(X_i - \bar{X})^2}$. Since $\mathrm{E}[Y_i] = \beta_0 + \beta_1 X_i$, we get $\mathrm{E}[Y_i - \bar{Y}] = \beta_1(X_i - \bar{X})$. Then $\mathrm{E}[\hat{\beta}_1] = \frac{\beta_1 \sum(X_i - \bar{X})^2}{\sum(X_i - \bar{X})^2} = \beta_1$. Thus $\hat{\beta}_1$ is an unbiased estimator of $\beta_1$.

\item \textbf{Exercise 4.} Suppose $X \sim \mathrm{Poisson}(\lambda)$ with observed value $x = 5$ in a single observation. Find the maximum likelihood estimate of $\lambda$. \emph{Solution.}$L(\lambda) = \frac{\lambda^5 e^{-\lambda}}{5!}$. Then $\ell(\lambda) = 5\ln\lambda - \lambda - \ln(5!)$. Setting $\frac{d\ell}{d\lambda} = \frac{5}{\lambda} - 1 = 0$ gives $\hat{\lambda} = 5$. The MLE is the observed count.

\item \textbf{Exercise 5.} A 95\% confidence interval for a population mean $\mu$ based on $n = 25$ observations with $\bar{x} = 10$ and $s = 4$ is $(\bar{x} - t_{0.025, 24} \cdot s/\sqrt{n},\; \bar{x} + t_{0.025, 24} \cdot s/\sqrt{n})$. Using $t_{0.025, 24} \approx 2.064$, compute the interval. \emph{Solution.}$s/\sqrt{n} = 4/5 = 0.8$. The margin of error is $2.064 \times 0.8 = 1.6512$. The confidence interval is $(10 - 1.6512,\; 10 + 1.6512) = (8.35,\; 11.65)$.

\end{enumerate}
\section*{8. References}
\begin{thebibliography}{9}
\bibitem{statistics_ref1} D. R. Cox and D. V. Hinkley, \emph{Theoretical Statistics}, Chapman and Hall/CRC, 1974.
\bibitem{statistics_ref2} E. L. Lehmann and J. P. Romano, \emph{Testing Statistical Hypotheses}, Springer-Verlag, 2005.
\bibitem{statistics_ref3} W. G. Cochran, \emph{Sampling Techniques}, Wiley, 1977.
\bibitem{statistics_ref4} G. Casella and R. L. Berger, \emph{Statistical Inference}, Cengage Learning, 2002.
\bibitem{statistics_ref5} D. A. Freedman, \emph{Statistical Models: Theory and Practice}, Cambridge University Press, 2009.
\end{thebibliography}
